\chapter*{Resumen}
\addcontentsline{toc}{chapter}{Resumen}

Este Trabajo Fin de Máster aborda la creación de lenguajes de dominio
específico con sintaxis basada en bloques. Bibliotecas como Blockly facilitan la
creación de editores visuales, pero todavía que el desarrollador defina
manualmente los bloques, los campos, las conexiones, la paleta, las
validaciones y los generadores en JavaScript. Este esfuerzo debe repetirse cada
vez que se necesita un editor para un nuevo dominio.

Para reducir este esfuerzo, se propone Model2Blockly, un entorno generativo
distribuido como plugin de Eclipse y como cadena de generación web. La
herramienta admite dos rutas de entrada: metamodelos Ecore enriquecidos con
anotaciones y descripciones textuales \texttt{.m2b} definidas con Xtext. La
ruta es un DSL de definición de lenguajes para describir qué bloques,
campos, conexiones y validaciones debe tener el editor generado. Ambas rutas se
normalizan en un metamodelo EMF intermedio, \texttt{EditorSpec}, que describe
bloques, campos, categorías, entradas de valor, referencias, validaciones y
metadatos de interfaz. Sus instancias se serializan como XMI y se recargan como
contrato entre la transformación de modelos y la generación final. A partir de
ese modelo se generan editores HTML/JavaScript que pueden abrirse en el navegador y
permiten guardar, cargar, exportar JSON y XMI de instancia, mostrar avisos de
validación e inspeccionar las reglas como bloques.

La solución se evalúa mediante un ejemplo conductor, AppMaker, descrito tanto
como metamodelo Ecore anotado como mediante el DSL textual \texttt{.m2b}. Este
dominio cubre pantallas, widgets, acciones, navegación, fuentes de datos,
contenciones, herencia, clases abstractas, referencias dinámicas, campos
obligatorios, validaciones simples y plantillas de código. La memoria compara
ambas formas de autoría y muestra cómo convergen en el mismo modelo
intermedio antes de generar el editor Blockly. El proyecto incluye un update
site p2 público, informes de generación, editores
desplegados en GitHub Pages y 142 pruebas JUnit organizadas por capas. Los
resultados indican que Model2Blockly permite generar editores Blockly a partir
de descripciones de dominio de alto nivel, reduce la programación manual y
mantiene una arquitectura MDE reutilizable.

\noindent\textbf{Palabras clave:} lenguajes de dominio específico, programación
basada en bloques, Blockly, EMF, Ecore, Xtext, ingeniería dirigida por modelos.

\chapter*{Abstract}
\addcontentsline{toc}{chapter}{Abstract}

\begin{english}
This Master's Thesis addresses the creation of domain-specific languages with a
block-based syntax. Libraries such as Blockly help developers build visual
editors, but they still require developers to define blocks, fields,
connections, toolboxes, validations and generators manually in JavaScript. This
work must be repeated whenever an editor is needed for a new domain.

To reduce this effort, the thesis proposes Model2Blockly, a generative
environment distributed as an Eclipse plugin and as a web generation pipeline.
The tool supports two input routes: annotated Ecore metamodels and textual
\texttt{.m2b} descriptions defined with Xtext. The textual route is a language-definition
DSL for specifying the blocks, fields, connections and validations of the
generated editor. Both routes are normalised into an intermediate EMF
metamodel, \texttt{EditorSpec}, which describes blocks, fields, categories,
value inputs, references, validations and interface metadata. Its instances are
serialized as XMI and reloaded as the contract between model transformation and
final generation. From this model, the tool generates HTML/JavaScript editors
that can be opened in a browser and support
saving, loading, JSON and domain instance XMI export, validation warnings, and
inspection of validation rules as blocks.

The solution is evaluated through a running example, AppMaker, described both
as an annotated Ecore metamodel and as a textual \texttt{.m2b} language
definition. This domain covers screens, widgets, actions, navigation, data
sources, containments, inheritance, abstract classes, dynamic references,
required fields, simple validations and code templates. The thesis compares the
two authoring forms and shows how they converge on the same intermediate model
before generating the Blockly editor. In addition, the project includes a
public p2 update site, generation reports, editors deployed through GitHub Pages,
and 142 JUnit tests organised by architectural layer. The results indicate that
Model2Blockly can generate Blockly editors from high-level domain descriptions,
reduces manual programming and keeps a reusable MDE architecture.

\noindent\textbf{Keywords:} domain-specific languages, block-based programming,
Blockly, EMF, Ecore, Xtext, model-driven engineering.
\end{english}
